% =====================================================================
%  Sec. 8 -- Physical validation: the two gaps scale in opposite directions
%  Replaces the "Validation" subsection of the former paper/results.tex
%  (lines 1-46 plus the four spectral-figure captions' physics claims).
%
%  NEW \label's introduced here (none exist in the v2 tree):
%      sec:physics, tab:gaps, tab:poles
%
%  RESOLVED 2026-09-18: all five entries below now exist in the bibliography, with
%      titles pulled from Crossref rather than guessed.  Two corrections came out of
%      that check and are applied in the prose of this section:
%        benthien2007 is the EXTENDED Hubbard model, computes THREE channels
%          (S(q,w), N(q,w), A(q,w)), and runs at L = 30 to 120, not "L=90-120";
%        senechal2000 states NO sixteen-site ceiling -- it works at 12 sites and argues
%          that cluster perturbation theory removes the need for a large cluster.  It is
%          cited in Sec. 8.6 for spin-charge separation, which it does establish, and the
%          size convention now rests on xu2024 alone.
%        \cite{xu2024}       Xu, Lu, Tohyama, Shao, arXiv:2407.13136 (2024)
%        \cite{senechal2000} Senechal, Perez, Pioro-Ladriere, PRL 84, 522 (2000)
%        \cite{benthien2007} Benthien & Jeckelmann, PRB 75, 205128 (2007),
%                            DOI 10.1103/PhysRevB.75.205128
%        \cite{essler2005}   Essler, Frahm, Gohmann, Klumper, Korepin, CUP 2005,
%                            DOI 10.1017/CBO9780511534843
%        \cite{kim1996}      C. Kim et al., PRL 77, 4054 (1996)
%      Existing keys used: liebwu1968, jeckelmann2002, nocera2016,
%      descloizeaux1962, mueller1981, haydock1972, golub2010, damascelli2003.
%
%  PROVENANCE.  Charge block of Table~\ref{tab:gaps}: release/data/charge_gap.json
%  (checks_status ALL PASS).  Spin and Heisenberg blocks, and every fit quoted in
%  Sec. 8.1: p2_calc/gap_scaling.json (generated 2026-09-18T11:45:32Z, full
%  reorthogonalization, self-checks pass).  Pole census: p2_calc/pole_structure.json
%  (generated 2026-09-18T15:05:07Z, n_l=120, dense anchors at L=6 and L=8).
%  DEPOSIT, checked 2026-09-18: DISCHARGED.  release/data/gap_scaling.json,
%  release/data/pole_structure.json and release/data/charge_gap.json are present, with
%  gap_scaling.py, pole_structure.py and charge_gap_ed.py in release/src/.  I re-read
%  pole_structure.json while editing this section and it reproduces Table~\ref{tab:poles}
%  entry for entry: 7 momenta 3-7 for A(k,omega), 6 momenta 6-9 for charge, 6 momenta 1-3
%  for spin.
%
%  FLOATS rebuilt in parallel, NOT written here: \ref{fig:lattice},
%      \ref{fig:sqw}, \ref{fig:spin} (its present bold title must go; see 8.2),
%      \ref{fig:method} (panel (a) caption replacement follows Sec. 8.4 below,
%      commented out at the end of this file), \ref{fig:akwsampled}.
%      *** If the figure agent renders the gap scaling as a panel, it must be
%      drawn as POLE MARKERS, not as a broadened spectrum.  Sec. 8.1 survives
%      Jeckelmann's rule only because it reads pole positions off the
%      tridiagonal recursion, which are independent of eta; a broadened-lineshape
%      panel would forfeit that defence entirely. ***
%
%  GREP NOTE for the mechanical audit: the three withdrawn terms are named
%  exactly ONCE each in this file, inside quotation marks, in the retraction
%  sentence of Sec. 8.2. A retraction cannot be written without naming what it
%  retracts. Nowhere else in this file is any of them used assertively.
% =====================================================================

\section{Physical validation: the two gaps scale in opposite directions}\label{sec:physics}

The two collective channels of the half-filled Hubbard chain are the strongest physical check
available to this method, because their answer is known exactly and independently of it. This section
reports that check as a scaling statement across $L=6,8,10,12$ rather than as a picture at a single
size, does not use three words that a single size cannot support, and
separates what is a validation of the method from what is a result. On the same lattice and from the
same primitive, the charge scale converges to a finite thermodynamic value while the spin scale closes
as $\approx2/L$, with an independent effective-Heisenberg control (Table~\ref{tab:gaps}).

\subsection{Charge converges to a finite value, spin closes as \texorpdfstring{$1/L$}{1/L}}\label{sec:gaps}

\begin{table*}[tp]\centering\small
\caption{\textbf{The two collective scales, across four system sizes.} Charge: the Mott gap
$\Delta=\mu^{+}-\mu^{-}$ from exact $(N{\pm}1)$ ground-state energies (half-filled Hubbard ring,
$U/t=8$; particle--hole symmetry $\mu^{+}+\mu^{-}=U$ verified to $10^{-8}$ at every size, ARPACK
checked against dense \texttt{eigh} on every sector of dimension $\le5000$). $\Delta_\infty=4.67952\,t$
is the Lieb--Wu thermodynamic value~\cite{liebwu1968}, evaluated here from its exact integral
representation to an absolute quadrature error of $2\times10^{-9}$ and cross-checked against a second
exact representation, the two agreeing to $10^{-13}$. Spin: the lowest pole of
$S^{zz}(q=\pi,\omega)$, a pole position of the tridiagonal recursion and therefore independent of the
broadening used to display it. Right-hand block: the same quantity in the pure Heisenberg ring at the
effective exchange $J=4t^{2}/U=0.5\,t$, a model with no charge degree of freedom at all. The residual
column is quoted to four decimals so that it is the difference of the two columns to its left; at five
it would not be.}\label{tab:gaps}
\begin{tabular}{lcccccc}
\toprule
 & \multicolumn{2}{c}{charge} & \multicolumn{2}{c}{spin (Hubbard)} & \multicolumn{2}{c}{spin (Heisenberg, $J=0.5\,t$)}\\
\cmidrule(lr){2-3}\cmidrule(lr){4-5}\cmidrule(lr){6-7}
$L$ & $\Delta/t$ & $(\Delta-\Delta_\infty)/t$ & $\Delta_{s}/t$ & $L\,\Delta_{s}/t$ & $\Delta_{\rm ST}/t$ & $L\,\Delta_{\rm ST}/t$\\
\midrule
$6$  & $5.35813$ & $0.6786$ & $0.34857$ & $2.091$ & $0.34237$ & $2.054$\\
$8$  & $5.18584$ & $0.5063$ & $0.23958$ & $1.917$ & $0.26134$ & $2.091$\\
$10$ & $5.02224$ & $0.3427$ & $0.19956$ & $1.996$ & $0.21162$ & $2.116$\\
$12$ & $4.96876$ & $0.2892$ & $0.16635$ & $1.996$ & $0.17792$ & $2.135$\\
\bottomrule
\end{tabular}
\end{table*}

\paragraph{Charge.} $\Delta(L)$ decreases monotonically and stays above the Lieb--Wu value
$4.67952\,t$ at every size, the residual falling from $0.679\,t$ at $L=6$ to $0.289\,t$ at $L=12$,
$6.2\%$ of the thermodynamic value; $\Delta(L=12)=4.96876\,t$ is the Hubbard-band separation of
Fig.~\ref{fig:lattice}. No extrapolation exponent is fitted: a straight line in $1/L$ undershoots
Lieb--Wu by $0.12\,t$, which is evidence that the form is not $1/L$, not evidence about the limit.

\paragraph{Spin.} The lowest pole of $S^{zz}(q=\pi,\omega)$ gives $L\,\Delta_{s}=2.091$, $1.917$,
$1.996$, $1.996$: a $1/L$ closing with a coefficient of about $2$ (mean $2.000$, range
$1.917$--$2.091$), and a fit $L\,\Delta_{s}=d_0+d_1L$ has a slope consistent with zero
($d_1=-0.010\pm0.018$). A linear-in-$1/L$ intercept of $\Delta_s$ is $-0.020\pm0.020$, against
$4.561\pm0.044$ for the charge gap; but the pure Heisenberg ring, whose gap is known to vanish, gives
$0.0142\pm0.0014\,t$ at these sizes, so a vanishing intercept is not by itself evidence that a gap
closes, and the flatness of $L\,\Delta_s$ is what carries the statement. That Heisenberg ring at
$J=4t^{2}/U=0.5\,t$---a different Hamiltonian, Hilbert space and code path---reproduces the $1/L$ law,
$L\,\Delta_{\rm ST}=2.054$--$2.135$, differing from the Hubbard value by $6.5\%$ at $L=12$ with a
residual that changes sign: it confirms the scaling and the coefficient to within $10\%$ and does not
pin the coefficient (Sec.~\ref{sm:physics}).

The two channels therefore behave \emph{qualitatively oppositely} under a change of system size---one
converging to a finite thermodynamic value, the other extrapolating to zero as $1/L$---with a factor
$29.9$ between them at $L=12$, and with $1/L$ intercepts that separate. That
opposition, measured at four sizes with an independent control on the smaller scale, is the physical
statement this work makes. It is a statement about pole positions, not about lineshapes, so it does
not depend on the broadening at which either response is plotted.

\subsection{What a single size at \texorpdfstring{$\eta=0.10\,t$}{eta = 0.10 t} could not support}

A spin--charge statement cannot be read off a lineshape at a single size, for a reason of resolution. At $L=12$ the spin scale is $0.16635\,t$ against the $\eta=0.10\,t$ used for
Fig.~\ref{fig:spin}: the gap is $1.66\,\eta$, i.e.\ $0.83$ of the Lorentzian full width at half
maximum. At that resolution a spin response with a $0.17\,t$ gap and one with none are not
distinguishable by eye in a single-size plot. The figure could \emph{illustrate} the physics; it could
not establish it. Three terms---``continuum'', ``gapless'' and ``edge of the continuum''---are therefore not licensed by it, and none is used for any channel of this work at this size.

What establishes the physics instead is the four-size scaling of Table~\ref{tab:gaps}, plotted in
Fig.~\ref{fig:gapscaling}, and a pole census (Table~\ref{tab:poles}): between $L=6$ and $L=12$ the
sector grows by a factor $2440$ while the largest per-momentum count of poles above $1\%$ of the channel
maximum grows from $3$ to $7$ in the addition branch, from $4$ to $9$ in the charge channel and from $2$
to $3$ in the spin channel, a discrete spectrum whose lines are broadened together rather than a dense
one forming. The census, the sum rule as measured, and the size convention of the
exact-diagonalization literature together with the dynamical-DMRG resolution
criterion~\cite{jeckelmann2002} this work does not meet at $L=12$, are in Sec.~\ref{sm:physics}.

\subsection{Spin--charge separation is a check of the method, not a result}

That the spin and charge responses of the one-dimensional Hubbard model live on separate energy scales
has been known exactly since Lieb and Wu~\cite{liebwu1968}, is textbook~\cite{essler2005}, was seen in
angle-resolved photoemission thirty years ago~\cite{kim1996,damascelli2003}, is already legible in the
spectral weight of twelve-site clusters embedded by cluster perturbation theory~\cite{senechal2000},
and has been computed by dynamical DMRG since 2007~\cite{benthien2007} in all three of the
momentum-resolved channels reported here---in that reference's notation the charge factor
$N(q,\omega)$ at $L=120$, the spin factor $S(q,\omega)$ at $L=100$ and the one-particle function
$A(q,\omega)$ at $L=90$, which are respectively our $S(q,\omega)$, our $S^{zz}(q,\omega)$ and our
$A(k,\omega)$, for the extended Hubbard model of which ours is the $V=0$ line; the same literature treats the Hubbard chain at $L=48$ with $\eta=0.1\,t$ by
correction-vector DMRG, and reaches $L=64$ in the Heisenberg chain~\cite{nocera2016}. We present it here as what it is: a check of the exact full-sector references, against which the sampled reconstructions are scored, against a known answer, with the des Cloizeaux--Pearson lower
boundary~\cite{descloizeaux1962} and the M\"uller upper bound~\cite{mueller1981} recovered as
thermodynamic-limit references rather than as findings. No part of it is offered as new physics.

% =====================================================================
%  DROP-IN CAPTION REPLACEMENT for panel (a) of \ref{fig:method}
%  (results.tex lines 44 and 51).  Not part of Sec. 8; move into the figure file.
%  It removes BOTH prohibited phrases that survive in that caption today:
%  "to machine precision" (line 44) and "complete-subspace" (line 51).
%
% \textbf{(a)} Sum rule. The analytic value $\int A^{+}d\omega=1-\langle\hat n_p\rangle=0.5$ is an
% identity of the ground state and is not a check of the reconstruction. What is plotted is the
% integral of the two curves over the $600$-point window at $\eta=0.15\,t$: $0.486588$ for the exact
% reference and $0.486471$ for the reconstruction, a relative difference of $2.4\times10^{-4}$. Both
% lie $2.7\%$ below $0.5$ because the Lorentzian tails leave the window; the trapezoidal residue is
% $3\times10^{-6}$ relative (Sec.~\ref{sec:physics}). The reconstruction shown uses $245$ of the $300$
% determinants of the sector, $82\%$, and reaches relative $L_1=0.017$.
%
%  NOTE ON THE TWO SUBSPACE SIZES IN THAT CAPTION: 245 and 242.5 are DIFFERENT objects and both are
%  correct.  245 (81.7%) is reconstructed_subspace_size in data/method_max.json, the panel-(a)
%  subspace, rel-L1 = 0.01727 -> the caption's 0.017.  242.5 is the 8-seed mean sampled |S| at K=12
%  in the same file, rel-L1 = 0.018550 -> the caption's 0.019.  An earlier internal audit treated
%  242.5 as the panel-(a) size; it is not.  Neither is "complete": the sector has 300 determinants.
% =====================================================================

